\documentclass[11pt,a4paper]{article}
\newcommand{\fichetype}{Exercices}
% Thème A4 commun, librement modifiable.
\usepackage[T1]{fontenc}
\usepackage[utf8]{inputenc}
\usepackage[french]{babel}
\usepackage{lmodern,amsmath,amssymb,mathtools,tikz,pgfplots}
\usepackage[margin=2.2cm,headheight=15pt]{geometry}
\usepackage{xcolor,fancyhdr,enumitem,hyperref,needspace}
\usetikzlibrary{arrows.meta,positioning,calc}
\pgfplotsset{compat=1.18,width=11cm,height=6cm}
\definecolor{Ink}{HTML}{202124}
\definecolor{Blue}{HTML}{356B9E}
\definecolor{Muted}{HTML}{737579}
\color{Ink}
\renewcommand{\familydefault}{\sfdefault}
\setlength{\parindent}{0pt}
\setlength{\parskip}{6pt}
\setlist{itemsep=3pt,topsep=4pt}
\pagestyle{fancy}\fancyhf{}
\fancyhead[L]{\small\color{Muted}MATHS LYCÉE}
\fancyhead[R]{\small\color{Muted}2026--2027}
\fancyfoot[L]{\small\color{Muted}\fichetype}
\fancyfoot[R]{\small\thepage}
\renewcommand{\headrulewidth}{0.3pt}
\hypersetup{colorlinks=true,urlcolor=Blue,linkcolor=Blue}
\newcommand{\R}{\mathbb{R}}
\newcommand{\N}{\mathbb{N}}
\newcommand{\E}{\mathbb{E}}
\newcommand{\V}{\mathrm{V}}
\newcommand{\exercice}[3]{\Needspace{6\baselineskip}\vspace{0.5cm}{\color{Blue}\small\textbf{EXERCICE #1}\hfill #2}\par{\large\bfseries #3}\par}

\begin{document}
{\small\color{Blue}PREMIÈRE SPÉCIALITÉ}\par{\LARGE\bfseries Droites et cercles dans un repère}\par{\large Exercices}\par\vspace{0.2cm}\hrule\vspace{0.4cm}
Les exercices progressent de l’application directe vers la recherche. Détailler les calculs et justifier les réponses.
\par\noindent\begin{minipage}{\linewidth}\exercice{1}{Application}{Lire et vérifier}
Pour $d:4x-3y+6=0$, donner une normale et un directeur. Les points $A(0,2)$ et $B(3,5)$ appartiennent-ils à $d$ ?
\end{minipage}\par
\par\noindent\begin{minipage}{\linewidth}\exercice{2}{Application}{Une droite par deux points}
Déterminer une équation de la droite passant par $A(-1,2)$ et $B(3,4)$. Donner une équation de sa parallèle passant par $C(0,-1)$.
\end{minipage}\par
\par\noindent\begin{minipage}{\linewidth}\exercice{3}{Entraînement}{Perpendiculaire et intersection}
Soit $d:2x-y+1=0$. Écrire une équation de la perpendiculaire à $d$ passant par $P(1,1)$, puis calculer leur intersection.
\end{minipage}\par
\par\noindent\begin{minipage}{\linewidth}\exercice{4}{Entraînement}{Projeté et distance}
Calculer le projeté $H$ de $P(4,0)$ sur $d:x+y-2=0$. Déduire la distance de $P$ à $d$ et justifier que tous les autres points de $d$ sont plus éloignés.
\end{minipage}\par
\par\noindent\begin{minipage}{\linewidth}\exercice{5}{Entraînement}{Médiatrice}
Déterminer la médiatrice de $A(-2,1)$ et $B(4,3)$. Vérifier que $M(0,5)$ appartient à cette droite en utilisant deux méthodes.
\end{minipage}\par
\par\noindent\begin{minipage}{\linewidth}\exercice{6}{Application}{Identifier des ensembles}
Identifier les ensembles définis par : (a) $x^2+y^2-6x+2y+6=0$ ; (b) $(x-1)^2+(y-2)^2=0$ ; (c) $(x+2)^2+y^2=-1$.
\end{minipage}\par
\par\noindent\begin{minipage}{\linewidth}\exercice{7}{Approfondissement}{Une route et une zone circulaire}
Une zone a pour frontière $(x-2)^2+(y-1)^2=25$. Une route suit $y=4$. Déterminer les deux points d’entrée et de sortie de la zone, puis la longueur de route située dans le disque.
\end{minipage}\par
\par\noindent\begin{minipage}{\linewidth}\exercice{8}{Approfondissement}{Tangente et vérification}
Pour le cercle $x^2+y^2=10$, vérifier que $A(1,3)$ est sur le cercle et trouver sa tangente. Montrer par substitution que la tangente n’a qu’un point commun avec le cercle.
\end{minipage}\par
\end{document}
